\section{Requiems, Ritual Masses, and Exceptional Actions}
\label{sec:requiem}

\subsection{Requiem classes}

All Souls Masses agree with the Office; every other Requiem is outside its order (rubric 390). No commemoration of the current Office is made in a Requiem (391). Every Requiem, \emph{even a funeral Mass}, is first subject to rubric 393: it is prohibited during exposition of the Blessed Sacrament and in the specified one-Mass conflicts with a required conventual Mass, the Mass following candle or ash blessing, or a required Rogation action. The class-specific rules below add to, rather than replace, that gate.

\begin{longtable}{>{\bfseries\raggedright\arraybackslash}p{0.14\linewidth}>{\raggedright\arraybackslash}p{0.28\linewidth}>{\raggedright\arraybackslash}p{0.42\linewidth}>{\raggedright\arraybackslash}p{0.08\linewidth}}
\toprule
Class & Included Mass & Admission rule & Rubrics\\
\midrule
\endhead
I & All Souls; one funeral Mass directly connected to a funeral & Subject to 393. All Souls follows its day. Funeral Mass then follows the specific rules in 406--407, not the general votive table. & 402--409\\
II & Day of death; after notice; final burial & Subject to 393; must be applied for that deceased; prohibited on every I-class day and every Sunday. & 410--414\\
III & Third, seventh, thirtieth day; anniversary; cemetery; within the eight-day period counted inclusively from All Souls & Subject to 393; prohibited on I- and II-class days; each subsection supplies its formularies and transfer possibilities. & 415--422\\
IV & Other daily Requiems & Subject to 393; IV-class ferias only, outside Christmas time; should genuinely be applied for the dead. & 423\\
\bottomrule
\end{longtable}

\subsection{All Souls}

Each priest may celebrate three Masses using the three appointed formularies and the ordering in rubrics 403--404. A priest celebrating one uses the first; one celebrating two uses the first and second. The sung or conventual Mass uses the first, with the permissions for ordering the others stated in the rubric. If All Souls occurs on Sunday, rubric 96b moves it to Monday as its proper seat. This is a transfer of the liturgical day, not an external solemnity or a private Requiem.

\subsection{The funeral Mass}

One Requiem directly joined to the funeral is I class (402, 405). Rubric 406 prohibits it:

\begin{itemize}
\item on precedence-table positions 1--6;
\item on holy days of obligation included among the other universal I-class feasts at position 11;
\item on the dedication anniversary and title of the church where the funeral occurs;
\item on the principal patron of the city or place;
\item on the title and canonized founder of the religious institute to which the church belongs.
\end{itemize}

Rubric 407 adds two easily missed qualifications. A funeral Mass is prohibited on a Sunday when an external solemnity of a feast is celebrated. If a feast named in rubric 406 is accidentally transferred, the funeral Mass remains prohibited on the feast's impeded original day but is permitted, other bars aside, on the day to which the feast's Office transfers.

An ordinary II-class Sunday is not by itself in rubric 406's enumeration. Textually, therefore, the 1960 Code can admit a funeral Mass on such a Sunday only when no rubric-393 bar, rubric-407 external solemnity, or named local or obligatory feast applies; it prohibits the Mass on an Advent, Lenten, Passiontide, Easter, Low, or Pentecost Sunday. Because this result is pastorally surprising and present discipline is outside this manual, an actual use must be checked against the competent Ordo and authority. When prohibited or reasonably impossible at the funeral, rubric 408 permits transfer to the nearest similarly unimpeded day.

\subsection{Other Requiems}

A death-day or final-burial Mass of II class is excluded by any Sunday or I-class day (411). III-class anniversaries and the third, seventh, and thirtieth days require a III- or IV-class day; rubrics 417 and 419 provide their own transfer options. A IV-class daily Requiem is narrower still and cannot be used in Christmas time (423).

One cannot reason, “The funeral Requiem is I class, therefore all I-class rules apply to every Mass for this person.” Each Requiem kind has a definition, application requirement, prohibition list, and formulary.

\subsection{Requiem assembly}

\begin{itemize}
\item Omit Psalm 42 and the Introit doxology (425b, 428).
\item Omit Gloria and Creed (432d, 476f).
\item Make no commemoration of the Office (391); Requiem prayers are governed by 398.
\item \latin{Dies irae} is obligatory only in I-class Requiems; it may be omitted in II--IV (399).
\item Use the Requiem Preface (499).
\item Follow the Requiem Agnus Dei and peace variants printed in the Order.
\item Say \latin{Requiescant in pace}; omit the blessing (507c, 508).
\item Omit the Last Gospel when the absolution at the bier follows; otherwise apply 509--510 exactly.
\end{itemize}

\subsection{Nuptial Mass}

The Nuptial Mass is a II-class votive when admitted, but it is prohibited on every Sunday (341, 379). If the nuptial blessing is permitted while the Mass itself is impeded, celebrate the occurring Mass, join the Nuptial Collect under one conclusion, and give the appointed blessings (378--381). On All Souls and during the Sacred Triduum, the votive, its commemoration, and the blessing within Mass are all prohibited. The repository's separate Nuptial Mass guide provides the complete verified ritual sequence; this manual uses it as an example of why a ritual formulary cannot be reduced to ten propers.

\subsection{Rogations, exposition, and other ritual settings}

Where the Major or Minor Litanies have the procession or ordered supplications, the Rogation Mass is a II-class votive and forms part of the whole action (346--347). Eucharistic exposition has distinct rules depending on Forty Hours, a full day, a few hours, and the altar of exposition (348--355). Consecration, blessing of a church or altar, ordination, profession, and grave public cause likewise have enumerated rules. Always identify:

\begin{enumerate}[leftmargin=1.45em]
\item the action and its competent authorization;
\item whether the Mass is inseparable, required, merely permitted, or replaced by the occurring Mass;
\item its votive class;
\item any ritual Collect joined under one conclusion;
\item the insertion point of every prayer or blessing;
\item prohibitions caused by the day or a one-Mass obligation.
\end{enumerate}

\subsection{Churches with only one Mass}

Rubrics 304, 326, and 393 protect the conventual Mass and certain blessings, Litanies, and actions where a church has only one Mass. A Mass category that is permitted in general can therefore be unavailable at that altar because the sole Mass must fulfill another obligation. “Permitted somewhere on this class of day” is not yet “permitted as this church's only Mass.”
